    Este trabalho investiga se traços de personalidade do modelo Big Five,
    induzidos por instrução (\textit{prompt}) em agentes de inteligência
    artificial baseados em modelos de linguagem de grande escala (LLM),
    predizem o perfil de investidor desses agentes, medido pelo instrumento
    IPA (Índice de Perfil do Investidor). A pergunta é desdobrada em três
    questões: se os LLMs simulam consistentemente os traços solicitados,
    quais associações entre personalidade e perfil de investidor são
    robustas e em que medida elas dependem do modelo e dos parâmetros de
    geração. Na análise principal, foram gerados 300 agentes com o modelo
    \texttt{llama3.1:8b}, cada um com níveis (Baixo/Médio/Alto) sorteados
    para os cinco traços. Os agentes responderam ao BFI-44 (autorrelato de
    personalidade, usado como checagem de manipulação) e ao IPA. A análise
    foi repetida com outras temperaturas de geração, com um novo sorteio de
    perfis e com outros modelos de linguagem. A indução de personalidade foi
    fiel nos cinco traços ($\rho$ entre 0,813 e 0,929). Na análise
    principal, Abertura à Experiência, Extroversão, Consciência e
    Neuroticismo predisseram significativamente o escore do IPA (regressão
    OLS, $R^2=0{,}459$; $p<0{,}001$), mas apenas Abertura à Experiência e
    Extroversão se mantiveram significativas em todas as rodadas da análise
    de robustez (temperaturas, sorteio e modelos); os efeitos de Consciência e
    Neuroticismo mostraram-se instáveis, e Amabilidade não teve efeito
    consistente. O perfil de investidor absoluto atribuído a uma mesma
    personalidade variou fortemente entre modelos. O achado de Abertura à
    Experiência converge com De Bortoli et al. (2019), Jiang, Peng e Yan
    (2024) e estudos recentes com LLMs. Conclui-se que agentes de IA
    baseados em LLM, quando adequadamente induzidos e com a manipulação
    verificada, reproduzem a direção de parte das relações conhecidas entre
    personalidade e tolerância ao risco, mas que o nível de risco que
    declaram é, em grande parte, uma característica do modelo usado.

    \vspace{\baselineskip}

    \noindent
    \textbf{Palavras-chave}: Big Five. Perfil de investidor. Inteligência
    artificial. Modelos de linguagem. Finanças comportamentais.
